\section{Indicator inequalities}\label{sec:moments}
\begin{lemma}[Positive-measure inequalities]\label{core:lem:measure-tools}
Let $\mu$ be a finite positive measure on the unit circle, with mass
$M$ and real first two moments $A,B$. Then
\[
 0\le M+B\le2M,\qquad 2A^2\le M(M+B).
\]
If its first $m$ moments all equal the real number $c$, then
\[
 \mu(\{1\})\le\frac{M+mc}{m+1}.
\]
These statements also apply after subtracting any known positive
atomic submeasure.
\end{lemma}
\begin{proof}
Writing $z=e^{it}$, one has
$\int\cos^2t\,d\mu=(M+B)/2$; Cauchy--Schwarz gives the first
inequalities. The nonnegative Fej\'er polynomial
$F_m(z)=|1+z+\cdots+z^m|^2/(m+1)$ has value $m+1$ at one and
integral $M+mc$, proving the atom bound.
\end{proof}

\begin{lemma}[Second moment]\label{core:lem:family-signed-second}
Let $u,v$ be the ordinary and signed inverse-square codegree sums.
For any signed dimension character,
\[
 E(u+v-1)\le\sum_{i=0}^3|\beta_i|
 \le1+\sqrt{3(E-1)}.
\]
\end{lemma}
\begin{proof}
The eight constituents of $I(1)$ have positive dimensions $E/f$
and squared mass $H=E^2(u+v)$. Their translates by the canonical
central sign object have opposite dimensions, are disjoint from
them, and have twists $\pm1$. The center's total squared dimension
is $4E^2$. The central indicator formula
\cite[Theorem~4.1]{NS} gives
\[
 \sum_A (\dim A)^2\theta_A^2=4Es,
 \qquad s=1+\sum_{i=1}^3\nu_2(X_i^+)\beta_i.
\]
Here each $\nu_2(X_i^+)=\pm1$ and the two lifts have opposite
dimensions and opposite second indicators. Removing the sixteen
known objects, the weighted sum of squared real parts of the
remaining twists is
\[
 \tfrac12\bigl((4E^2-2H)+(4Es-2H)\bigr)
 =2E^2+2Es-2H\ge0.
\]
This gives the first bound, and Cauchy's inequality gives the second.
\end{proof}
\begin{lemma}[Quartic trace and norm]\label{core:lem:family-trace-norm}
If $L$ has degree four, the signed rational algebra is $L$ itself,
and its four character codegrees are the conjugates of $E$, with
multiplicity. Moreover
\[
 \Tr_{L/\Q}(E)\le\sum_j f_j,
 \qquad \Norm_{L/\Q}(E/f)\in\ZZ_{>0}
\]
for every ordinary codegree $f$. If $v$ is the signed inverse-square
sum, then
\[
 v\ge\frac12-\frac92\,
 \frac{\Tr_{L/\Q}(E)}{\Norm_{L/\Q}(E)}.
\]
\end{lemma}
\begin{proof}
The dimension character is a surjection from a rational algebra of
dimension four onto $L$, hence an isomorphism. The Casimir's regular
trace is the sum of the codegrees. Its trace is also
$\sum_i\Tr(T_i^2)=\sum_{i,j,k}(T_{ij}^k)^2$, at most the
corresponding ordinary sum. Norm integrality is
Lemma~\ref{core:lem:family-induction}. For $q_i=1/g_i$ the reciprocal
identity gives $\sum q_i=1$, and the elementary four-variable
Schur inequality $9e_3(q)+2\sum q_i^2\ge1$ gives the result,
since $e_3(q)=\Tr(E)/\Norm(E)$.
\end{proof}

\subsection{Galois averages in a quartic signed field}
Assume here that the signed dimension field $L$ has degree four.
It identifies the signed rational algebra with $L$. We use the
normalized trace identity~\eqref{core:frobenius-trace}.
\begin{proposition}\label{core:prop:quartic-individual-disjointness}
Suppose that, for every integer $1\le n\le m$,
\begin{equation}\label{core:eq:quartic-individual-intersection}
 L\cap\Q(\zeta_n)=\Q.
\end{equation}
Then
\begin{equation}\label{core:eq:quartic-averaged-bound}
 2u+v\le\frac34+\frac3{m+1}.
\end{equation}
\end{proposition}

\begin{proof}
Choose a finite Galois extension $\Omega/\Q$ containing $L$,
all twists of simple objects of $Z(\mathcal S)$, and
$\Q(\zeta_n)$ for $1\le n\le m$.  Put
$G=\operatorname{Gal}(\Omega/\Q)$.  On the unit circle define
\[
 \mu=\frac4{|G|}\sum_{\sigma\in G}
       \sum_{A\in\operatorname{Irr}Z(\mathcal S)}
       \frac{\sigma(\dim A)^2}{4\sigma(E)^2}\,
       \operatorname{Dirac}_{\sigma(\theta_A)}.
\]
Every central dimension lies in $L$, since forgetting expresses it as
an integral linear combination of the $x_i$.  All its conjugates are
real, so this measure is positive.  The global-dimension identity for
the center shows that its total mass is four.

The central-induction indicator formula \cite[Theorem~4.1]{NS} is
$\operatorname{tr}(\theta^n|_{I(V)})=2E\nu_n(V)$.
Multiplying by $\dim V$ and summing, adjunction gives
\[
 \sum_A\dim(A)^2\theta_A^n
 =2E\sum_{V\in\operatorname{Irr}\mathcal S}
                    \dim(V)\nu_n(V).
\]
Consequently, for $1\le n\le m$,
\begin{equation}\label{core:eq:quartic-moment-indicator}
 \int z^n\,d\mu(z)
 =\frac2{|G|}\sum_{\sigma\in G}
       \sum_{V\in\operatorname{Irr}\mathcal S}
       \sigma\!\left(\frac{\dim V}{E}\right)
       \sigma\bigl(\nu_n(V)\bigr).
\end{equation}
The cyclic operator defining $\nu_n(V)$ has order dividing $n$;
hence $\nu_n(V)\in\Q(\zeta_n)$
\cite[Theorem~5.1]{NSHigher}.  By
\eqref{core:eq:quartic-individual-intersection}, restriction of $G$ to
$L$ and $\Q(\zeta_n)$ is onto the product of their Galois
groups.  The average of each product in
\eqref{core:eq:quartic-moment-indicator} therefore factors into the two
separate averages.  Equation~\eqref{core:frobenius-trace}
annihilates all nonunit terms.  The two remaining objects are
$\mathbf1$ and $\varepsilon$.  The unit has every indicator equal to
one.  The sign has indicator $-1$ at even orders and zero at odd orders.
Their normalized dimension traces are respectively $1$ and $-1$.
It follows that
\begin{equation}\label{core:eq:quartic-constant-moments}
 \int z^n\,d\mu(z)=
 \begin{cases}1,&n\text{ even},\\[2pt]\tfrac12,&n\text{ odd}.
 \end{cases}
\end{equation}
Lemma~\ref{core:lem:translate-twists} gives the indicator signs,
the eight constituents of $I(1)$ of dimensions $E/f_i,E/g_i$,
and their eight disjoint sign translates.
Thus translates of the $R$ constituents have twist $+1$, while
translates of the $T$ constituents have twist $-1$.  These statements
are preserved by every $\sigma\in G$.  Summing the squared weights
yields the two estimates
\begin{equation}\label{core:eq:quartic-atom-bound}
 \mu(\{1\})\ge 2u+v,\qquad \mu(\{-1\})\ge v.
\end{equation}

By Cauchy--Schwarz, $v\ge1/4$.  Consequently
$\nu=\mu-\tfrac14\operatorname{Dirac}_{-1}$ is a positive measure.
Its mass is $15/4$, and its first $m$ moments are all $3/4$.
Apply Lemma~\ref{core:lem:measure-tools} to obtain
$(m+1)\mu(\{1\})\le(15+3m)/4$. Together with
\eqref{core:eq:quartic-atom-bound}, this proves the claim.
\end{proof}


\begin{proposition}[Squared moments]\label{core:squared-moments}
If $L\cap\Q(\zeta_{2j})=\Q$ for $1\le j\le a$, then
\[
u+v\le\frac12+\frac{3}{2(a+1)}.
\]
In particular $u+v\le7/8$ always, and $u+v\le4/5$ if
$\sqrt2\notin L$.
\end{proposition}
\begin{proof}
Push the preceding averaged measure forward by $z\mapsto z^2$.
Its total mass is four, its first $a$ moments are one, and its atom
at one is at least $2(u+v)$. Apply Lemma~\ref{core:lem:measure-tools}. The first three
intersections are trivial because $L$ is totally real. The fourth
can be nontrivial only if $\Q(\sqrt2)\subset L$.
\end{proof}
\begin{corollary}\label{core:quartic-short-moments}
If the signed dimension field $L$ is quartic, then
\[
 \sqrt2\notin L\ \Longrightarrow\ u+v\le\frac45,
 \qquad
 \sqrt5\notin L\ \Longrightarrow\ 2u+v\le\frac98.
\]
\end{corollary}
\begin{proof}
The first assertion is Proposition~\ref{core:squared-moments}.
For the second, the only nontrivial totally real intersection of $L$
with a cyclotomic field of order at most seven can be
$\Q(\sqrt5)$. Apply
Proposition~\ref{core:prop:quartic-individual-disjointness} with $m=7$.
\end{proof}

\begin{lemma}\label{core:lem:sqrt5-periodic-moments}
Suppose the signed dimension field is $K=\Q(\sqrt2,\sqrt5)$.
With the preceding notation, one has
\[
 5u+4v\le\frac{15}{4}.
\]
\end{lemma}
\begin{proof}
Use the unsquared averaged central measure $\mu$ above.
The nontrivial quadratic subfields of $K$ have conductors $5,8,40$.
Hence $K\cap\Q(\zeta_n)=\Q$ whenever $5\nmid n$ and
$8\nmid n$. The same indicator calculation gives
\[
 \int z^n\,d\mu=\begin{cases}1/2&n\text{ odd},\\1&n\text{ even},\end{cases}
 \qquad 5\nmid n,\quad8\nmid n.
\]
These identities hold at every such order: extending the averaging
field leaves the measure unchanged.

Restrict $\mu$ to the fortieth roots of unity. Averaging the functions
$z^{40j}$ over $j$ gives their indicator function, since the measure
has finite root-of-unity support. Thus, for each residue $n$ divisible
by neither $5$ nor $8$, the restricted measure retains the displayed
$n$th Fourier coefficient. Subtract
$\nu=\tfrac34\operatorname{Dirac}_1+\tfrac14\operatorname{Dirac}_{-1}$.
The resulting forty weights have Fourier support only at multiples
of $5$ or $8$. By the discrete Fourier transform they are therefore
\[
 r_{ij}=A_i+B_j,\qquad i\in\ZZ/8,\quad j\in\ZZ/5,
\]
where a fortieth root with exponent $l$ corresponds to
$(i,j)=(l\bmod8,l\bmod5)$.
All cells other than $(0,0)$ and $(4,0)$ are nonnegative. The forced
atoms give
\[
 r_{00}\ge2u+v-\frac34,\qquad
 r_{40}\ge v-\frac14.
\]
The total of the forty weights is at most $4-1=3$.
For an additive array there is the exact identity
\[
 \sum_{i,j}r_{ij}-5r_{00}-3r_{40}
 =\frac12\sum_{j\ne0}r_{4j}
  +\frac54\sum_{\substack{i\ne0,4\\j\ne0}}r_{ij}\ge0.
\]
Consequently
$3\ge5(2u+v-3/4)+3(v-1/4)$, proving the assertion.
\end{proof}
\subsection{Pivotal positivity and second-indicator parity}
\begin{lemma}[A pivotal Fej\'er bound]\label{core:lem:family-pivotal-fejer}
For a pseudo-unitary self-dual rank-four category,
\[
 (9u-3)D\le6+2\sqrt{3(D-1)}.
\]
The same holds for the FP data whenever a spherical dimension
character is Galois conjugate to the FP character.
\end{lemma}
\begin{proof}
Let $A_j$ be the constituents of $I(1)$, with dimensions $D/f_j$.
Their known twist-one mass $L_X$ in $I(X)$ satisfies
\[
 \sum_{X\ne1}d_X L_X=\sum_j(D/f_j)(D/f_j-1)=D^2u-D.
\]
Apply $|1+z+z^2|^2$ to the positive twist weights of $I(X)$.
The first indicator vanishes, the second is $\pm1$, and the
polynomial has value nine at one. Thus
$9L_X\le D(3d_X+2)$. Multiplying by $d_X$ and summing gives
$9(Du-1)\le3(D-1)+2\sum_{X\ne1}d_X$; Cauchy's inequality
finishes. For a conjugate dimension character, apply the embedding
to the induction and indicator identities. Forgetful multiplicities
are nonnegative integers, so all embedded central dimensions are
positive; twists remain roots of unity and second indicators remain
signs. The same proof applies.
\end{proof}


We will also use the unsimplified form of the same three-term estimate:
\begin{equation}\label{core:three-term-unsimplified}
9(D^2u-D)\le D\bigl(3(D-1)+2S\bigr),
\qquad S=\sum_{X\ne1}d_X.
\end{equation}
It follows directly from the proof of
Lemma~\ref{core:lem:family-pivotal-fejer} before applying Cauchy's inequality.
\subsection{Quadratic comparison and determinant concentration}
\begin{lemma}[Quadratic comparison]\label{core:lem:family-comparison}
Let $D,D'$ be the ordinary FP orbit and $E\ge E'$ a proper signed
dimension orbit in the same quadratic field. Set
$P=D^{-1}+D'^{-1}$, $q=E^{-1}+E'^{-1}$, $N=DD'$ and
$k=EE'/N$. Then $k$ is a positive integer,
\begin{equation}\label{core:quadratic-strict}
 kq<P+(k-1)/D,\qquad q\le2P/3.
\end{equation}
\end{lemma}
\begin{proof}
Norm divisibility gives $k\in\ZZ_{>0}$. Since
$D>E\ge\sqrt{kN}$, compare $x+kN/x$ at $E$ and $D$ to obtain
the strict inequality. If $k\ge2$, then $D/D'>k$ and
$1/D<P/(k+1)$, whence $q<2P/(k+1)\le2P/3$.
If $k=1$, $s=D/E>1$ is a totally positive unit, so its integral
trace is at least three. Using $E\ge E'$ gives
$q/P\le2/(s+s^{-1})\le2/3$.
\end{proof}


